\documentclass[12pt,a4paper,openany]{ctexbook}
\usepackage{geometry}\geometry{margin=1in}
\usepackage{hyperref}
\usepackage{setspace}\onehalfspacing
\usepackage{fancyhdr}\pagestyle{fancy}
\fancyhf{}\fancyhead[LE,RO]{\thepage}\fancyhead[RE]{\leftmark}
\renewcommand{\headrulewidth}{0.4pt}
\usepackage{graphicx}
\usepackage{xcolor}
\usepackage{enumitem}
\definecolor{veilgold}{RGB}{180,140,60}
\definecolor{veildark}{RGB}{30,20,40}

\begin{document}

\title{帷幕之地 · UI原画需求规格书}
\author{帷幕之地开发组}
\date{\today}
\maketitle
\tableofcontents
\newpage

\chapter{总体美术风格}

\section{风格定位}

帷幕之地的美术风格定位为\textbf{暗黑奇幻 × 东方玄幻}的融合体。它既不是纯粹的西方暗黑哥特，也不是传统的东方仙侠——而是这两种风格在“末世之后”语境下的有机融合。

\textbf{核心关键词：}撕裂感、残存、灵焰微光、暗而不脏、古而不旧。

\section{色调体系}

\begin{itemize}
    \item \textbf{主色调：}深紫黑（夜空）+ 暗金（灵焰微光）+ 冷灰（骨灰平原）
    \item \textbf{辅色调：}血月红（蚀月之光）+ 霜青（永冬山脉遗存）+ 墨绿（幽翠残林）
    \item \textbf{高光色：}淡金色（纯净灵焰）+ 银白（帷幕纤维）
    \item \textbf{暗部色：}近黑紫（蚀痕边缘）+ 深蓝黑（无光之渊）
\end{itemize}

\section{光照原则}

所有画面中的光源必须是有\textbf{叙事来源}的——不是凭空的环境光，而是画面中某个具体事物发出的光。主要光源类型：

\begin{enumerate}
    \item \textbf{灵焰光}——人物体内发出的淡金色/淡蓝色微光，柔和，从内部向外透
    \item \textbf{血月光}——天空中的红色月光，冷而锐利，拉长所有阴影
    \item \textbf{帷幕极光}——天边的带状流动光幕，深紫到墨绿渐变
    \item \textbf{炉火光}——庇护所内的暖橙色火光，唯一让人感觉安全的颜色
    \item \textbf{蚀痕黑光}——蚀者身上的“反向光”——不是照亮周围，而是吸收周围的光
\end{enumerate}

\chapter{角色原画：八位隐藏职业}

\section{通用要求}

\begin{itemize}
    \item 每个角色至少需要一张\textbf{半身立绘}（用于角色卡显示）
    \item 每个角色建议追加一张\textbf{全身动态图}（用于夜晚行动界面）
    \item 画布尺寸：半身立绘 1200×1600px，全身动态 1600×2400px
    \item 所有角色需要有明确的\textbf{灵焰可视表现}——在其胸口/心脏位置，有一团微弱的、半透明的火焰光芒透出衣物
    \item 纯净者的灵焰为淡金色/暖白色/淡蓝色（因人而异）；蚀者的灵焰为纯黑色带暗紫边缘，且表现为“吸收光而非发出光”
\end{itemize}

\section{蚀者 CORRUPTED}

\textbf{角色定位：}灵焰被彻底侵蚀的凡人，必须吞噬他人灵焰为生。狼人杀原型的继承者——但不再是狼形怪物，而是被黑暗从内部吞噬的人类。

\textbf{外貌特征：}
\begin{itemize}
    \item 曾经是普通人——蚀者保留了他们蚀变前的外貌，所以每个蚀者看起来都不一样
    \item 但在灵视中——他们的灵焰是纯黑的，像一口微型黑洞，把周围的光向内吸收
    \item 瞳孔边缘有一圈难以察觉的暗紫色微光——只有在彻底黑暗的环境中才能看到
    \item 皮肤颜色比正常人浅一层——因为灵焰不再供给身体温暖，皮下血管隐约透出一种淡灰调
    \item 不是“变丑”——是“变空了”。像一座有人住过但现在已经搬空了的房子
\end{itemize}

\textbf{原画要求：}
\begin{itemize}
    \item 画一个曾经是普通村民、现在已蚀变的人。不画“怪物”——画一个你会在聚落中擦肩而过而不会多看一眼的人。但如果你盯着看——你会觉得什么东西不对劲
    \item 他的灵焰（胸口位置）是一片不反射任何光的黑色区域。它不发光——它“吸光”。周围空气中的微尘似乎都在向那片黑色微微倾斜
    \item 面色平和——甚至可能带着一丝微笑。但眼睛里没有笑意
\end{itemize}

\section{冥僧人 NETHER MONK}

\textbf{角色定位：}曾侍奉诸神的僧侣，在神死后与帷幕之外的古老存在做了交易。以残存人性碎片换取堕化他人之力。一个“知道自己正在变成什么、但选择继续借贷”的悲剧存在。

\textbf{外貌特征：}
\begin{itemize}
    \item 穿着一件曾经可能是僧袍但现在已经被各种来源不明的污渍、烧痕、和不属于这个世界的物质浸染成灰黑的长袍
    \item 脖子上挂着一串由不同材料（木珠、骨片、金属碎屑——每颗代表一次借贷）串成的念珠。念珠中的某些珠子在黑暗中偶尔闪烁
    \item 眼睛——这是他最显著的特征。一只眼睛正常（人眼），另一只眼睛变成了完全的暗紫色——没有瞳孔，没有巩膜，就是一面平滑的暗紫色表面。那是在第一次与幕外存在做交易时被“签约”的印记
    \item 手持一根长杖——杖头不是装饰，而是一团被锁在透明琥珀中的、缓慢蠕动的蚀痕样本
\end{itemize}

\textbf{原画要求：}
\begin{itemize}
    \item 半身立绘：画出冥僧人平静但令人不安的面容。他的正常眼睛看着你，他的签约眼睛看着你身后的某个地方——也许是帷幕之外。嘴角微抿——不是愤怒，不是悲伤，是“我知道你不知道的事情”
    \item 长杖可见——杖头琥珀中的蚀痕应该被画成一种被压扁的、不断缓慢变换形状的黑色液滴
    \item 念珠中的某几颗发出极其微弱的、不同颜色的光——代表他还没还清的几次借贷
\end{itemize}

\section{帷幕学者 VEIL SCHOLAR}

\textbf{角色定位：}研究帷幕运作规律的人。能通过“察灵”感知他人灵焰是否被蚀痕污染。继承了奈莉娅观测塔的学术传承。

\textbf{外貌特征：}
\begin{itemize}
    \item 穿着朴素的学者袍——深灰或深棕色，材质粗糙但干净。领口处绣着一只微型眼睛符号——帷幕学者的标志
    \item 戴着单片眼镜——不是装饰，镜片是由观测塔废墟中找到的霜星水晶磨制的。透过镜片看到的灵焰是彩色的（普通人只能看到淡金）。镜框是一种暗银色的灵质金属
    \item 其中一只眼睛的瞳孔边缘有些微不规则——那是长期凝视裂隙导致的水晶体微损伤。但这不是弱点——这只眼睛获得了看到裂隙波动的能力
    \item 随身携带一本厚重的笔记——皮质封面，里面的纸页参差不齐（来自不同年代的补页）。封面上用暗金线绣着“帷幕观测录”几个字
\end{itemize}

\textbf{原画要求：}
\begin{itemize}
    \item 半身立绘：学者正低头看着笔记，单片眼镜中倒映出一小片淡金色的灵焰光芒。背景可以是观测塔的内部——风化的石墙、被风吹乱的纸页、窗口外是深紫色的帷幕天空
    \item 一只眼睛要能看到微小的灵焰残像——在镜片的那一侧，空中应该悬浮着一些半透明的、淡金色的、正在消逝的残迹
\end{itemize}

\section{草药学者 HERBAL SAGE}

\textbf{角色定位：}研究帷幕之地灵植的学者。能炼制蚀灭符（摧毁蚀者）和愈灵符（修复灵焰）。原毒巫的能力继承者——但不再是煮毒药的女巫，而是绘制灵符的学者。

\textbf{外貌特征：}
\begin{itemize}
    \item 穿着带有多个口袋的实用外套——每个口袋装着不同的灵植标本、符纸、或研磨工具。外套的肘部和膝盖处有补丁——是长期跪在地上采集灵植留下的磨损
    \item 手指上有被灵植汁液染出的淡淡的颜色——浅绿、淡紫、微黄——洗不掉的，是职业印记
    \item 腰间挂着一个皮制符囊——里面存放着预先绘制好的灵符。符囊在黑暗中会从缝隙中透出微弱的金色或暗紫色光（愈灵符是金色，蚀灭符是暗紫色）
    \item 手持一支符笔——笔杆是世界树残枝，笔头是灵质冰原狼的尾毛。正在空白的符纸上绘制符文
\end{itemize}

\textbf{原画要求：}
\begin{itemize}
    \item 半身立绘：草药学者蹲在幽翠残林的边缘，一只手刚摘下一株发着微光的灵植，另一只手正在空白的符纸上绘制。符纸上已经能看到一个正在形成的金色符文——那种光芒不刺眼，是柔和的、像日出前东方第一缕光
    \item 周围地面上应该散落着几片用过的旧符——符上的灵光已经熄灭，但符纸上的符文还在，可以被后来者拾取学习
\end{itemize}

\section{愈灵师 SPIRIT MENDER}

\textbf{角色定位：}能够直接触碰他人灵焰的稀有天赋者。能修复受损灵焰，也能通过蚀痕净化仪式摧毁蚀者。原药巫的能力继承者——但不再是“用药”的人，而是直接以自身灵焰作用于他人灵焰的人。

\textbf{外貌特征：}
\begin{itemize}
    \item 外观普通到你可能在聚落中见过十次也记不住——愈灵师不需要引人注目，他们的工作从来不是被看到的
    \item 但他们的双手是不同的——长期直接触碰他人灵焰使他们的手掌皮肤呈现出一种半透明感。在黑暗中，能看到他们掌心的血管中流淌着一种微微发光的、淡金色的液体——那是残留的被愈者的灵焰碎片
    \item 不戴手套——因为任何物理阻隔都会干扰灵焰触觉的精度。双手是他们唯一的工具
    \item 眼睛总是微微眯着的——不是在审视，是在“感觉”。愈灵师的灵焰感知是持续激活的，他们时刻知道周围人的灵焰状态
\end{itemize}

\textbf{原画要求：}
\begin{itemize}
    \item 半身立绘：愈灵师正伸出手，手掌向上，掌心微微发光。从手掌中延伸出极细的、几乎看不到的金色丝线——这些丝线正伸向画面外的某个受伤者的灵焰。愈灵师的表情是专注的——不是紧张，是那种“我已经做过几千次了但仍然不敢大意”的专注
    \item 背景可以暗示灵焰的存在——不是画火，是画空气中有一种微微的热浪般的折射
\end{itemize}

\section{帷幕守卫 VEIL GUARDIAN}

\textbf{角色定位：}曾在帷幕裂隙边缘服役的战士。能将自身灵焰延伸出体外形成保护屏障。以自身灵焰的磨损为代价来守护他人。

\textbf{外貌特征：}
\begin{itemize}
    \item 穿着一套由锻甲灵织者打造的灵质铠甲——不是全身板甲，而是覆盖关键部位（胸、肩、前臂）的轻甲。甲片上有微小到肉眼几乎看不到的、反复折叠锻造留下的纹路——那是灵质金属的“记忆”
    \item 左手持一面中型盾——盾面是骨灰平原上挖出的冥界天秤碎片铸成的。盾面上能看到微弱到只有在特定角度才可见的、安努赫卡的审判符文
    \item 右臂有旧伤——手臂上有一些已经愈合但永远不会消失的裂纹状的疤痕。那不是物理伤害——那是他的灵焰多次撕裂后在身体上留下的映射
    \item 站姿——不是高傲的英雄站姿，而是一个随时准备挡在前面的人。重心微微前倾，盾牌始终在自己和目标之间
\end{itemize}

\textbf{原画要求：}
\begin{itemize}
    \item 半身立绘：帷幕守卫挡在画面最前方，盾在前，他的灵焰从胸口延伸出来——形成一层淡金色半透明的、覆盖在盾面上的额外屏障。屏障上有细微的裂缝——那是之前守护留下的损伤
    \item 背景是夜晚的聚落——隐约可以看到身后有一间庇护所，窗户中透出微弱的炉火光
\end{itemize}

\section{灵痕追猎者 FLAME TRACKER}

\textbf{角色定位：}能感知灵焰被夺走时残留痕迹的猎手。使用灵焰猎枪（远程）和噬灭短铳（近身反击）。原猎人的继承者——枪械保留，但弹药变成了灵焰。

\textbf{外貌特征：}
\begin{itemize}
    \item 穿着适合在荒原和残林中行动的猎装——深色，耐磨，有很多口袋装弹药和追踪工具。披着一件由灵质冰原狼皮制成的短斗篷——白色，在月光下会微微反光
    \item \textbf{灵焰猎枪}——这是一把长管步枪，枪身是黑森林枯死圣橡的心材雕刻的，枪管是矿脉山脉的灵质金属锻造的。枪托上刻着被这把枪击杀的蚀者的标记。枪不需要子弹——发射的是追猎者自身灵焰的微小碎片。扣动扳机时，枪管中会短暂地闪过一道淡金色的光
    \item \textbf{噬灭短铳}——挂在腰间的短管火铳，由骨灰平原深处挖出的冥界天秤碎块铸造。发射的不是弹丸，而是一束灵焰脉冲——在极近距离下，脉冲能震碎蚀者的蚀痕附着。只能使用一次（因为天秤碎片会暂时被蚀痕覆盖，需要愈灵师净化后才能再次使用）
    \item 眼睛——追猎者的眼睛在黑暗中能看到灵痕，所以在他们的视角中世界是被发光的丝线交织的。眼睛本身看起来和普通人一样，但瞳孔在黑暗中的反应比常人快得多
\end{itemize}

\textbf{原画要求：}
\begin{itemize}
    \item 半身立绘：追猎者半蹲在地面上，一只手按着地面——在感受昨晚留下的灵痕。另一只手握着灵焰猎枪，枪口朝向地面但随时准备抬起。他们的目光不是在看——是在“追踪”。眼神穿过画面，落在你身后的某个看不见的痕迹上
    \item 地面上应该有若隐若现的灵痕——一些淡金色或银白色的丝线，从追猎者的手边延伸到画面外。其中一根丝线是深黑色的——那是蚀者留下的
\end{itemize}

\section{灵织者 SPIRIT WEAVER}

\textbf{角色定位：}守幕者中数量最多的群体。拥有灵焰感知、帷幕低语捕捉、灵织术整合信息的能力。原村民的继承者——但不再是“普通人”，而是拥有独特灵焰感知天赋的聚落居民。

\textbf{外貌特征：}
\begin{itemize}
    \item 灵织者的外貌差异极大——因为每一个人都是一个独立的个体，不是职业的“模板”。但在所有灵织者身上——当他们激活灵焰感知时——他们的瞳孔中会出现极细的、像织机上的金线一样的反光
    \item 随身带着灵织工具——最典型的是一台可以挂在腰间的小型灵质织机，由黑森林枯枝和灵质丝线制成。织机不用于织布——它用于将感知到的灵焰碎片“编织”成可读的图案
    \item 八种子类型各有不同特征：
    \begin{itemize}
        \item \textbf{老兵灵织者}——年迈，身体有旧伤，但眼睛比任何人都锐利
        \item \textbf{行路灵织者}——腿上有长途跋涉留下的旧伤疤，背着旅行的行囊
        \item \textbf{学徒灵织者}——年轻，腰间挂着由愈灵师祝福过的药草包
        \item \textbf{记述灵织者}——随身带着纸笔，手指沾着墨水的痕迹
        \item \textbf{守夜灵织者}——眼窝深陷但从不闭上——他们已经忘了怎么完全睡着
        \item \textbf{炊火灵织者}——衣服上有面粉和灶灰的痕迹，手心常年是暖的
        \item \textbf{锻甲灵织者}——双臂粗壮，皮肤上有反复被灵质金属灼伤的斑痕
        \item \textbf{织幕灵织者}——最罕见——眼睛是双色的（一金一银），那是同时看到普通世界和灵质维度的标记
    \end{itemize}
\end{itemize}

\textbf{原画要求：}
\begin{itemize}
    \item 灵织者的原画可以画一组而非一个人——三四个不同类型的灵织者站在一起，在夜晚的聚落中，每个人都在做自己的事：一个蹲在地上感知灵焰波动，一个在记录信息，一个紧锁眉头在听帷幕低语。他们是在工作——不是在摆姿势。画他们工作时最自然的状态
\end{itemize}

\chapter{角色原画：十五位魂印者（表层身份）}

\section{通用要求}

\begin{itemize}
    \item 每位魂印者需要一张\textbf{半身肖像}（用于选择身份界面）
    \item 画布尺寸：1000×1200px
    \item 每位魂印者需要有明确的\textbf{魂印标记}——在他们的灵焰中（胸口透出的微光中），有一个独特的、只有该角色拥有的微小符号
    \item 背景——每位魂印者来自帷幕之域的不同区域（永冬山脉、幽翠残林、骨灰平原等），背景应暗示他们的出身地
\end{itemize}

\section{霜脊谱系}

\subsection{西格德 SIGURD —— 最后的霜脊猎手}
\begin{itemize}
    \item 老年男性，须发尽白，脸上有风霜刻下的深纹
    \item 穿着厚重的永冬皮草——由冰原狼皮和冰熊皮拼接的猎装
    \item 手持一把传承了七代的燧发猎枪——枪托上刻着七代猎人的名字，最新一个名字还很新
    \item 魂印标记：一个微小的冰晶六角形——代表他在冰原上追踪灵质冰原狼时冻结在灵焰中的那一刻时间
    \item 背景：永冬山脉的冰原，远处有一尊被冻结在奔跑姿势中的霜巨灵冰雕
\end{itemize}

\subsection{芙蕾维娅 FREYVIA —— 华纳遗孤}
\begin{itemize}
    \item 年轻女性，容貌中有一种不属于霜脊的柔和——那是已消失的麦野圣地的血统
    \item 头发是淡金色的——像被遗忘的麦浪的颜色。但这个金色正在逐渐变淡
    \item 灵焰是双色的——中心是麦野的金黄，外缘是霜脊的淡青——两种颜色在她的胸口处融合但不混合
    \item 魂印标记：一片金色的麦叶——那是她被从麦野圣地抛出时嘴里含着的最后一片故乡
    \item 背景：霜脊山脉南麓山坡上的草甸——她每年夏至都会来，试图回忆一个已经不存在的地方
\end{itemize}

\subsection{哈尔瓦德 HALVARD —— 裂隙之子}
\begin{itemize}
    \item 成年男性，生来就带有蚀痕——世界上唯一一个天生携带蚀痕还保持纯净的人类
    \item 脖子上戴着一条银链——链坠是一块只有指甲盖大的银片，内侧刻着“记住你是一个人”
    \item 银链在哈尔瓦德灵焰波动时会发光——不是在警告，是在“锁住”
    \item 魂印标记：一个被银链环绕的微型黑色球体——球体是蚀痕，银链是他的意志
    \item 背景：暮色聚落边缘——面朝帷幕站立，银链在血月下泛着微弱但坚定的银光
\end{itemize}

\section{翠庭谱系}

\subsection{莫莉安 MORIGHAN —— 最后的德鲁伊}
\begin{itemize}
    \item 年轻女性，绿色调的眼睛——在黑暗中会反射出极微弱的绿光
    \item 头发中编入了枯死圣橡的细枝——不是装饰，是她与幽灵树保持联系的“天线”
    \item 肩膀上和手臂上停着几只深色的鸟——不是渡鸦（莫丽甘娜的渡鸦早已灭绝），是一种在残林中新出现的、灰色的、叫声像哭声的鸟
    \item 魂印标记：一棵微型的圣橡树形——树的根系向下延伸，但枝叶是半透明的——因为那是灵质层面的枝叶
    \item 背景：幽翠残林深处——枯死的圣橡树干在月光下泛着银灰色的轮廓，空中悬浮着不可见的幽灵树的枝叶
\end{itemize}

\subsection{布丽吉德 BRIGID —— 圣火祭司}
\begin{itemize}
    \item 成年女性，红色长发——不是染的，是长期接触圣火后头发自然转成了暗红
    \item 手持一支从不熄灭的火把——燃料不是油，是希望。火把的火焰是暖橙色的，但在特定角度会变成白色
    \item 手心有被轻微灼伤的痕迹——不是因为不小心，是因为有时候圣火太旺了她必须用手盖住防止烧伤别人
    \item 魂印标记：一小簇始终在摇曳的火焰——不烫，不烧，只是不停地跳动着
    \item 背景：灰烬高地的圣火遗址——地面是焦黑的但还有余温，远处天边是帷幕的极光
\end{itemize}

\section{裂霆谱系}

\subsection{赫克托 HECTOR —— 特洛伊之盾}
\begin{itemize}
    \item 壮年男性，身形粗壮，脸上有一道从额头延伸到下巴的旧伤疤——那是年轻时为保护聚落被蚀者噬灵冲击留下的
    \item 左手持一面青铜盾——盾面上是一层淡金色的灵焰护盾——他的灵焰在盾面上延伸出保护屏障
    \item 站姿——不是高傲的战士站姿，而是一个不断在“等待下一击”的护卫。重心在脚尖，盾始终在前
    \item 魂印标记：一面微小的、圆形的、金边盾牌——盾面上有一道裂缝，但他从未让裂缝扩大
    \item 背景：暮色聚落的入口——他挡在聚落和无尽荒原之间
\end{itemize}

\section{金沙谱系}

\subsection{卡赫特 KAHET —— 冥途引渡者}
\begin{itemize}
    \item 年轻男性，深色皮肤，眼睛是淡金色的——那是长期在骨灰平原上凝视亡灵灵焰后眼睛被染上的颜色
    \item 额头上有一个由祖母临终前亲手画上的天秤符号——用骨灰平原的沙土和灵焰残液的混合物画的，已经渗入了皮肤，无法洗掉
    \item 腰间挂着一台小型天平——不是装饰，他能用它称量亡者灵焰中最后未说完的话的重量
    \item 魂印标记：一台微型的天秤——一边放着一片纯白羽毛，另一边空着——等着下一次称量
    \item 背景：骨灰平原——地面上是细密的灰白沙土，夜空中有一颗特别亮的星（那是拉赫卡里昂太阳舟的最后一道余晖）
\end{itemize}

\section{白石谱系}

\subsection{罗慕路斯 ROMULUS —— 双重灵焰者}
\begin{itemize}
    \item 成年男性，外表看不出任何异常——这本身就是最不正常的地方
    \item 他的灵焰是双色的——一半血红（来自罗马诸神的“标记”），一半淡金（他自己的纯净灵焰）。两种颜色不融合——像油和水在一个容器中共存
    \item 眼睛——在特定的光照角度下，左眼瞳孔会短暂地染成暗红色。他自己不知道——没有人告诉过他
    \item 魂印标记：两道互相缠绕但不交融的灵焰——一道血红，一道淡金——永远在旋转
    \item 背景：白石帝国遗址的律法广场——巨大的白色大理石地面，一个人站在曾经站过上万人的讲坛前
\end{itemize}

\section{永冬谱系}

\subsection{斯卡蒂 SKADI —— 雪山猎手}
\begin{itemize}
    \item 年轻女性，身形修长——在一个常年零下的环境中，瘦长的身体更容易在深雪中快速移动而不会沉下去
    \item 穿着纯白的猎装——由灵质冰原狼的皮制成。冰原狼的白色皮毛自带一种灵焰层面的伪装——身穿这种猎装的人，蚀者无法通过灵焰感知定位
    \item 身边站着一头灵质冰原狼——不是宠物，不是坐骑——是伙伴。冰原狼的灵焰是透明的白色，在雪地中完全不可见
    \item 魂印标记：一只冰原狼的爪印——不是刻上去的，是冰原狼第一次触碰她灵焰时留下的自然印记
    \item 背景：永冬山脉的冰原——极光在天空中展开它的翡翠和玫瑰色光幕。地面上有冰原狼的足迹，正在被新雪缓慢覆盖
\end{itemize}

\section{灰域谱系}

\subsection{伊瑟尔 YSERA —— 暗影织网者}
\begin{itemize}
    \item 年轻女性，头发是深色的，但在帷幕极光下会反射出暗紫色——那是长期在裂隙附近活动导致头发被灵质沁染
    \item 穿着不起眼的暗色调衣物——在聚落中你永远不会注意到她，这正是她想要的
    \item 腰间挂着一个小型织网——由灵质丝线编织成的网络，每一个节点代表着她收集到的一条秘密
    \item 魂印标记：一张微型的蛛网——网的每一根丝都在微微颤动，因为总有新的秘密落在上面
    \item 背景：帷幕裂隙边缘的废弃石屋——窗户外面可以看到流动的帷幕极光
\end{itemize}

\section{矿脉谱系}

\subsection{格伦 GOREN —— 灵质锻甲师}
\begin{itemize}
    \item 中年男性，双臂像树干一样粗壮——打了一辈子铁的人的手臂
    \item 皮肤上有多处被灵质金属灼伤的斑痕——那不是普通的烧伤。灵质金属在锻打时会在一定频率下发出灵质脉冲，长期暴露会使皮肤产生一种半透明的、像冰裂纹一样的永久性痕迹
    \item 背后是铁砧和锻炉——炉火中的铁块在发着微弱的灵质光
    \item 魂印标记：一把微型铁锤——锤头上有一道细到几乎看不见的裂缝，那是他用这把锤子锻造的第一件灵质门锁时——因为不熟练——锤子差点碎掉
    \item 背景：矿脉山脉下的铁匠铺——木架上钉着上百块被用旧的门锁，每一块后面都附着一小片皮纸
\end{itemize}

\section{骨灰谱系}

\subsection{艾琳 AILIN —— 守墓人}
\begin{itemize}
    \item 中年女性，穿着朴素到几乎可以融入骨灰平原沙土的衣物——灰色、棕色、土黄色
    \item 双手有长期挖掘留下的老茧——不是普通的挖土，而是用特殊工具在骨灰平原上挖掘灵焰残响“收容坑”
    \item 眼神中有一种说不上来的东西——不是悲伤（她已经悲伤了四十年，悲伤已经变成了她的底色），而是理解。她的眼睛看过太多次灵焰消散——所以她比任何人都更知道“活着”是什么
    \item 魂印标记：一把小铲子——交叉在一束已经枯萎但还没有完全化为灰的灵植上
    \item 背景：骨灰平原边缘——一排整齐的收容坑，每个坑旁边立着一块用旧了的石头，石头上刻着一个名字
\end{itemize}

\section{疾风谱系}

\subsection{奥里克 ORIC —— 疾风信使}
\begin{itemize}
    \item 年轻男性，身形瘦长但肌肉紧实——是一个将全身肌肉都优化为“奔跑”的人
    \item 腿上有多处旧伤——在荒原上奔跑时被瘴气、泥沼、荆棘留下的。不严重——因为他的速度够快，每次都在伤势加剧之前冲出了危险区域
    \item 背上绑着一捆密封的皮纸信件——每一封都是某个聚落居民托他送到另一个聚落的手写信。他没有打开过任何一封，但他知道每一封信放在哪个位置
    \item 魂印标记：两只翅膀——不是鸟的翅膀，是信使的飞翼靴的简化符号。两只翅膀是不同颜色的——一只是暮色聚落的土褐，另一只是远方聚落的淡灰
    \item 背景：平原驿站——一条被无数夜行者踩出的、浅浅的路径，夕阳（或者说日舟的最后一道余晖）正把它染成淡金色
\end{itemize}

\section{观测谱系}

\subsection{奈莉娅 NELIA —— 帷幕学者（始祖）}
\begin{itemize}
    \item 中年女性，一只眼睛是正常的深棕色，另一只眼睛——在一次近距离裂隙观测中被部分灼伤后——瞳孔变成了固定的半开状态，虹膜上环绕着三道暗紫色的环形疤痕
    \item 穿着实用的观测袍——多处被风吹旧了，袖口被墨水和灵质药水浸染出了洗不掉的痕迹
    \item 手中拿着她一直在写的那本《帷幕观测录》——封面上用金线绣的标题已经被磨得几乎看不清了
    \item 魂印标记：一只微型的眼睛——瞳孔中倒映着一道帷幕波纹。那是在她观测第三十一年的某个血月之夜，她终于看到了她想找的那个信号——千面之影的低频波动
    \item 背景：观测塔废墟——石塔朝向帷幕的那半面还完整，朝向聚落的那半面已经塌了。塔顶的木桌上散落着笔记、灵质水晶碎片、和一杯已经冷了的茶
\end{itemize}

\chapter{场景原画}

\section{暮色聚落全景}
\begin{itemize}
    \item \textbf{视角：}鸟瞰或半鸟瞰，从聚落边缘略高处俯视
    \item \textbf{画面元素：}十二栋屋子围成不规则的环形，中心是一棵枯死的黑色树干（世界树残枝），聚落外是荒原
    \item \textbf{氛围：}夜幕初降——天空是从暗蓝向血月红过渡的时刻。帷幕在地平线处发出微弱的、流动的极光。十二栋屋子中有几栋的窗户透出暖橙的炉火光——那是已经在守焰的人
    \item \textbf{用途：}大厅/等待界面背景
\end{itemize}

\section{观测塔}
\begin{itemize}
    \item \textbf{视角：}从塔底向上仰视，或塔内向外看
    \item \textbf{画面元素：}石塔——一半已塌、一半完好地指向天空。完好的那一半朝向远方的帷幕。塔顶的窗户透出微弱的观测光
    \item \textbf{氛围：}孤独但不阴森——一座被一个人住了三十一年的塔，充满了书、笔记、水晶碎片、空茶杯、和被风吹到每一个角落的纸张
    \item \textbf{用途：}夜晚序幕/学者相关场景
\end{itemize}

\section{聚落广场}
\begin{itemize}
    \item \textbf{视角：}人眼高度，站在广场中心
    \item \textbf{画面元素：}枯黑树干在正中，周围是十二栋屋子的模糊轮廓。地面上有被无数前人踩过的石径
    \item \textbf{氛围：}白昼——天空是淡灰蓝色，日舟的光透过帷幕后已经不那么明亮了，但依然温暖。聚落中的人在广场上走动——老人、孩子、灵织者、追猎者——所有人都在做自己的事
    \item \textbf{用途：}日间讨论/投票界面背景
\end{itemize}

\section{血月之夜}
\begin{itemize}
    \item \textbf{视角：}一个夜行者从庇护所的窗户向外看的视角
    \item \textbf{画面元素：}窗外是聚落的夜晚——月亮是深红色的（比普通满月大三倍），帷幕在地平线处发出强烈的、翻涌的极光。地面上可以看到几道若隐若现的发光灵痕——有的淡金（纯净者），有的深黑（蚀者）。远处有一间庇护所的窗户突然灭了
    \item \textbf{氛围：}紧张但不绝望——危险是真实的，但灵痕证明还有人活着、还在移动、还在守护别人
    \item \textbf{用途：}夜晚行动界面背景
\end{itemize}

\chapter{UI元素原画}

\section{角色卡片框架}
\begin{itemize}
    \item 每个角色的卡片需要一个统一的\textbf{框架纹理}——暗底色、金边，四个角上有类似于帷幕纤维的微小纹路
    \item 框架的上部是角色半身像的显示区域，下部是角色信息（名字、阵营、能力）
    \item 蚀者阵营的卡片框架微偏暗紫，守幕者阵营微偏暗金
\end{itemize}

\section{灵焰图标}
\begin{itemize}
    \item 一个通用的\textbf{灵焰图标}——用于表示角色的灵焰状态（纯净/被蚀痕侵蚀/正在修复）
    \item 纯净状态：淡金色，形状稳定，微微发光
    \item 蚀痕状态：黑色核心，暗紫色边缘，形状不稳定，像在吸收周围的光
    \item 修复状态：淡金色但有裂缝，裂缝中有愈灵的金色丝线在连接断口
\end{itemize}

\section{职业图标（8枚）}
\begin{itemize}
    \item 蚀者：一个被暗紫色环环绕的黑点——像一口正在吸入周围灵焰的微型裂隙
    \item 冥僧人：一只半闭的眼睛，瞳孔中有一个倒悬的、正在滴落的黑色液滴
    \item 帷幕学者：一只睁开的眼睛，瞳孔中有一个微小的帷幕波纹符号
    \item 草药学者：一片灵植的叶子——叶片上有微弱的金色脉络，叶尖处延伸为一道正在书写的符文
    \item 愈灵师：两只手掌心相对，中间有一团淡金色的灵焰——掌心的血管微微发光
    \item 帷幕守卫：一面圆形盾牌——盾面上有一层半透明的屏障，屏障上有微细的裂纹
    \item 灵痕追猎者：一支猎枪和一条向下延伸的金色灵痕——金色的线从枪口延伸到画面底边
    \item 灵织者：一台微型织机——经线上有不同颜色的灵质丝线（金、银、灰、黑）
\end{itemize}

\section{放逐裁决界面}
\begin{itemize}
    \item 聚落公投的界面——视觉上需要一个\textbf{裁决台}的意象
    \item 灵感来自骨灰平原上安努赫卡天秤的碎片——画面中可以有一个被风化了的、但仍然能看出天秤形状的石台
    \item 石台两侧各有一排等待投票的聚落居民——他们的轮廓是模糊的（因为他们还没有投票），只有已经投票的人轮廓会亮起来
\end{itemize}

\section{灵焰地图（小地图）}
\begin{itemize}
    \item 聚落的灵焰活动图——用于灵织者的灵焰感知界面
    \item 视觉上是深色底+发光线条——每条线代表一个人的灵焰移动轨迹。金色线是纯净者，灰色线是刚被蚀痕沾染的，黑色线是蚀者
\end{itemize}

\chapter{插图需求汇总}

\begin{center}
\begin{tabular}{|l|c|c|c|}
\hline
\textbf{类别} & \textbf{数量} & \textbf{尺寸(px)} & \textbf{优先级} \\
\hline
八职业半身立绘 & 8 & 1200×1600 & P0-必需 \\
八职业全身动态 & 8 & 1600×2400 & P1-建议 \\
十五魂印者半身肖像 & 15 & 1000×1200 & P0-必需 \\
暮色聚落全景 & 1 & 2400×1600 & P0-必需 \\
血月之夜晚景 & 1 & 2400×1600 & P0-必需 \\
聚落广场日景 & 1 & 2400×1600 & P1-建议 \\
观测塔场景 & 1 & 2400×1600 & P1-建议 \\
职业图标 & 8 & 128×128 & P0-必需 \\
灵焰状态图标 & 3 & 64×64 & P1-建议 \\
角色卡片框架 & 2 & 可平铺 & P0-必需 \\
放逐裁决背景 & 1 & 2400×1600 & P1-建议 \\
灵焰地图底图 & 1 & 800×600 & P2-可选 \\
\hline
\textbf{合计} & \textbf{50} & —— & —— \\
\hline
\end{tabular}
\end{center}

\chapter{交付格式}

\begin{itemize}
    \item \textbf{格式：}PNG（带透明通道——如果角色需要放在不同背景上）或 PSD（保留图层）
    \item \textbf{色彩空间：}sRGB
    \item \textbf{分辨率：}300dpi（打印级）或 72dpi（屏幕级）——建议同时提供两个版本
    \item \textbf{命名规范：}\texttt{category\_character\_variant.png}——例如 \texttt{role\_corrupted\_portrait.png}
\end{itemize}

\end{document}
